\documentclass[10pt,a4paper]{article}
\usepackage{amsmath,amssymb,amsthm,graphicx}
\usepackage{xfrac}
% Minimal editorial layout adapter for the pinned JEP source
% sources/symmetric-iet-infinite-jep-2025-302/source0.tex.
% The publisher document class (cedram) is not part of the retained material, and it is the
% class that used to define a handful of plain text abbreviations appearing in the author's
% prose. This adapter supplies exactly those abbreviations and nothing else: no mathematics,
% no notation, no theorem environment, no counter. The author's own macro block is retained
% separately from the pinned source itself.
% Corpus retention convention: CC BY 4.0, attribution to the original authors.
\providecommand{\ie}{i.e.\ }
\providecommand{\eg}{e.g.\ }
\providecommand{\resp}{resp.\ }
\providecommand{\skpt}{\smallskip}
\newenvironment{enumeratei}
{\bgroup\def\theenumi{\roman{enumi}}\def\theenumii{\arabic{enumii}}\begin{enumerate}}
{\end{enumerate}\egroup}

\hyphenation{addi-tion-al belong belongs denote exchan-ge exchan-ged exchan-ges}

\newcommand{\psfrac}[2]{\sfrac{(#1)}{#2}}

\newcommand\mto{\mathchoice{\longmapsto}{\mapsto}{\mapsto}{\mapsto}}
\let\oldbigsqcup\bigsqcup
\renewcommand{\bigsqcup}{\mathop{\textstyle\oldbigsqcup}\displaylimits}

\let\epsilon\varepsilon

\DeclareMathOperator{\Ess}{Ess}
\DeclareMathOperator{\Leb}{Leb}

\usepackage{tikz}
\usetikzlibrary{plotmarks}

\theoremstyle{plain}
\newtheorem{theorem}{Theorem}[section]
\newtheorem{lemma}[theorem]{Lemma}
\newtheorem{proposition}[theorem]{Proposition}
\newtheorem{corollary}[theorem]{Corollary}
\newtheorem{claim}{Claim}

\theoremstyle{definition}
\newtheorem{remark}[theorem]{Remark}
\newtheorem{example}[theorem]{Example}

\newcommand{\A}{\mathcal{A}}
\newcommand{\B}{\mathcal{B}}
\newcommand{\I}{\mathcal{I}}

\newcommand{\Z}{\mathbb{Z}}
\newcommand{\N}{\mathbb{N}}
\newcommand{\R}{\mathbb{R}}
\title{An effective criterion for ergodicity of the extension along a cocycle}
\author{Przemyslaw Berk, Frank Trujillo}
\date{doi:10.5802/jep.302; CC BY4.0}
\begin{document}
\maketitle
\noindent\textit{Editorial note.} The counted statement is the author's labelled proposition environment at source0.tex lines 1065-1082; its proof environment opens at line 1084, the line after the blank separator, so the pairing is by adjacency and the retained proof range 1084-1162 is the whole environment.\par\medskip
\noindent\textit{Editorial note.} Two theorem environments are retained as context because the counted statement and proof depend on them: thm: essvalues (source lines 1033-1035), which introduces the essential values, and thm: FUcriterion (source lines 1039-1050), whose four numbered conditions the counted statement invokes at source line 1069.\par\medskip
\noindent\textit{Editorial note.} This source has no prepared adapter and its preamble is built from the author's own macro and theorem block (source lines 8-45).\par\medskip
\noindent\textit{Editorial note.} This source has no label-number dump and the retained counter values are left empty, so the printed numbering of the card is the wrapper's default; the identity of the counted item is its author label.\par\medskip
\noindent\textit{Editorial note.} r002 note: r001 failed because the retained text uses the publisher class's plain text abbreviation resp (source line 901 in the proof, and source lines 575-576 in the retained prop:induced\_symmetric context), which the standalone wrapper does not define. r002 adds the minimal editorial adapter of this source as a first preamble span; the adapter contains only plain text abbreviations, no mathematics, no theorem environment and no counter, and nothing else in the card changes.\par\medskip
\noindent\textit{Editorial note.} r003 note: the r002 adapter span asked for line 13 of a twelve-line adapter file and the assembler refused it before compiling; r003 corrects the span to lines 1-12 and changes nothing else.\par\medskip
\noindent\textit{Editorial note.} r004 note: r003 failed with an undefined citation because the retained thm: FUcriterion context carries the author's citation of the Conze-Fraczek result; r004 bundles the publisher bibliography of the pinned source, which supplies exactly that entry, and changes nothing else.\par\medskip
\noindent\textit{Editorial scope.} One original result of the pinned publisher article, counted once: an effective ergodicity criterion for the extension of an interval exchange transformation along a differentiable cocycle (author label prop: effcriterion).\par\medskip
\noindent\textit{Allowed standard prerequisites (editorial).}\begin{itemize}
\item Skew product over an interval exchange transformation. The extension T\_f acts on the product of the interval with the torus by combining the base transformation with the cocycle f, and its ergodicity is measured with respect to the product of the invariant measures.. It is the transformation whose ergodicity is concluded.
\item Rokhlin towers with prescribed displacement. There are subsets that decompose into towers of bounded height, of small base and with uniformly small displacement over the tower height.. They are required by the second condition of the statement and are what makes the criterion effective.
\item Essential values. The essential values of the cocycle form a closed subgroup of the range, and ergodicity of the skew product is equivalent to that subgroup being the whole space.. It is the criterion that the effective conditions are designed to reach.
\end{itemize}\medskip
\begin{theorem}\label{thm: essvalues}
With the notation as above, the set $\Ess(T_f)$ is a closed subgroup of~$\R^m$. Moreover, $T_f$ is ergodic with respect to $\mu\otimes \Leb_{\R^m}$ if and only if $T$ is ergodic and $\Ess(T_f)=\R^m$.
\end{theorem}

\begin{theorem}[{\cite[Lem.\,2.7]{conze_cocycles_2011}}]\label{thm: FUcriterion}
Let $a\in\R^m$. Assume that for every $\epsilon>0$ there exists a sequence
of subsets $\{\Xi_n\}_{n\in\N}$ and an increasing sequence
$\{q_n\}_{n\in\N}$ of natural numbers such that
\begin{enumerate}
\item \label{cond:measure} $\liminf_{n\to\infty}\mu(\Xi_n)>0$,
\item \label{cond:rigidity}$\lim_{n\to\infty}\sup_{x\in \Xi_n}|T^{q_n}(x)-x|=0$,
\item \label{cond:almost_invariant} $\lim_{n\to\infty}\mu(\Xi_n\triangle T(\Xi_n))=0$,
\item \label{cond:BS} $|S_{q_n}f(x)-a|<\epsilon$.
\end{enumerate}
Then $a\in \Ess(T_f)$.
\end{theorem}

\begin{proposition}\label{prop: effcriterion}
Let $T:I\to I$ be an ergodic IET and let $m\in \N_+$. Assume that the function $f:I^{\times m}\to \R^{m}$ satisfies the following:
\begin{enumeratei}
\item \label{cond:continuity} $f$ is of the form $(x_1,\ldots,x_m)\mto (f_1(x_1),\ldots,f_m(x_m))$ with each $f_j$ being differentiable over the exchanged intervals,
\item \label{cond:sequences} there exist sequences $\{\Xi_n\}_{n\in\N}$ and $\{q_n\}_{n\in\N}$ satisfying \eqref{cond:measure}-\eqref{cond:almost_invariant} in Theorem \ref{thm: FUcriterion}, with the additional assumption that, for any $n \in \N$, $\Xi_n=\bigsqcup_{i=0}^{h_n-1}T^{i}({I_n})$ is a Rokhlin tower of $h_n\le q_n$ intervals with $|I_n|\le\sfrac{1}{q_n}$ and $\sup_{x\in \Xi_n}|T^{q_n}(x)-x|\le D/q_n$ for some $D>1$,
\item \label{cond:cont_interval} for every $x\in\Xi_n$, the interval $[x,T^{q_n}(x)]$ (or $[T^{q_n}(x),x]$) is a continuity interval of $f_j$, for every $j\in\{1,\ldots,m\}$,
\item \label{cond:derivative} there exists $C>1$ such that
\[
|S_{q_n}f_j(x)|\le C, \qquad C^{-1}q_n\le S_{q_n}f_j'(x)\le Cq_n,
\]
for any $j=1,\ldots,m$ and any $x\in\Xi_n$, and
\[
|f_j'(x)|\le C,
\]
for any $x\in I$.
\end{enumeratei}
If additionally $T^{\otimes m}$ is ergodic, then so is $T_f^{\otimes m}$.
\end{proposition}

\begin{proof}
In view of Theorem \ref{thm: essvalues}, it~is enough to show that we can
obtain an $m$\nobreakdash-dimensional cube of essential values. We will prove the
proposition for $m=1$ by realizing the essential values through Rokhlin
towers, which are subsets of $\Xi_n$. For $m\ge 2$ we first obtain the
same result for every $j=1,\ldots,m$, \ie we find an interval $(a_j,b_j)$
of essential values, realized through Rokhlin towers inside $\Xi_n$. Hence,
the set $(a_1,b_1)\times\ldots\times(a_m,b_m)$ is the set of essential
values realized via sequences of Rokhlin towers inside
$\bigsqcup_{i=0}^{h_n-1}(T^{\otimes m})^i(I_n^{\times m})$.

That being said, we~assume from now on that $m=1$. Note that due to our assumption on $\Xi_n$ and the derivative of $f$, each of the sets $S_{q_n}f(\Xi_n)\subseteq [-2C,2C]$ is a uniformly bounded union of intervals. Note that, since
\[
\liminf_{n\to\infty} h_n |I_n| = \liminf_{n\to\infty}\mu(\Xi_n) > 0
\]
and $h_n\le q_n$, we~have that, by~passing to a subsequence if necessary, there exists $E>1$ with
\[
E^{-1}\le q_n|I_n|\le 1.
\]
Hence, given the assumption on the derivative, we~have
\[
\Leb_{\R}\left(S_{q_n}f(\Xi_n)\right)\ge \frac{1}{CE}>0
\]
for every $n\in\N$.
By passing to a subsequence if necessary, we~may assume that the sets $S_{q_n}f(\Xi_n)$ have a nonempty intersection. Let $y\in \bigcap_{i=1}^{\infty} S_{q_n}f(\Xi_n)$. We claim that~$y$ is an essential value. For this purpose, we~will construct a subtower $\Xi_{n}^{y}\subseteq \Xi_n$, i.e., a~subset of $\Xi_n$ which itself is a Rokhlin tower of intervals, depending on $\epsilon>0$, such that \eqref{cond:measure} and \eqref{cond:BS} in Theorem \ref{thm: FUcriterion} is satisfied. This is enough, since \eqref{cond:rigidity} and \eqref{cond:almost_invariant} are automatically satisfied by any sequence of subtowers of $\Xi_n$.

Fix $\epsilon>0$ and consider the point {$x_n \in \Xi_n$ such that $S_{q_n}f(x_n)=y$}.
Let $\ell_n\in\{0,\ldots,h_n-1\}$ be such that $x_n\in T^{\ell_n}(I_n)$ and assume without loss of generality that $\ell_n<h_{n}/2$, the other case being treated symmetrically. Since each level of the tower is an interval then
\begin{align*}
\tag*{either}
\Bigl(x_n,x_n+\frac{1}{\max\{C, D,E\} q_n}\Bigr)\subseteq T^{\ell_n}(I_n),\\
\tag*{or} \Bigl(x_n-\frac{1}{\max\{C,D,E\} q_n},x_n\Bigr)\subseteq T^{\ell_n}(I_n).
\end{align*}
Again, without loss of generality, let us assume that it is the former. Consider the tower
\[
\Xi_n^y:=\bigsqcup_{i=0}^{\epsilon h_n/4CD} T^i\left(x_n,x_n+\frac{\epsilon}{2\max\{C,D,E\} q_n}\right)\subseteq \Xi_n.
\]
We now show that for every $x\in \Xi_n^y$ we have $S_{q_n}f(x)\in(y-\epsilon,y+\epsilon)$. If $x\in T^{\ell_n}(I_n)$, then by the mean value theorem, we~have
\[
\left|S_{q_n}f(x)-y\right|=\left|S_{q_n}f(x)-S_{q_n}f(x_n)\right|\le Cq_n|x-x_n|\le \epsilon/2.
\]
If $x\in T^{j} (T^{\ell_n}(I_n))$ with $j=1,\ldots,\epsilon h_n/4CD$, then we
split the Birkhoff sum into two pieces $S_{q_n}f(x) = S_{q_n - j}f (x) +
S_jf(T^{q_n - j}(x))$, which we estimate separately. For the first term, we~have
\[
\left|S_{q_n-j} f(x)-S_{q_n-j} f(T^j(x_n))\right|\le Cq_n|x-T^j(x_n)|\le\epsilon/2.
\]

For every $k\in\{0,\ldots,h_n/2\}$ we have that $T^{k}(x_n)$ and
$T^{k}(T^{-j}(x))$ belong to the same level of $\Xi_n$ and, since $\Xi_n$ is a
tower of intervals as such, belong to the same interval continuity of $f$.
Moreover, by~\eqref{cond:cont_interval}, $T^k(T^{q_n-j}(x))=T^k(T^{q_n}(T^{-j}(x)))$ and
$T^{k}(T^{-j}(x))$ belong to the same
continuity interval of $f$ for every $k\in\{0,\ldots,h_n/2\}$. Hence
$T^k(T^{q_n-j}(x))$ and $T^k(x_n)$ belong to the same continuity intervals of
$f$. Hence, by~\eqref{cond:BS} and mean value theorem we
get
\[
\left|S_j f(T^{q_n-j}(x))-S_j f(x_n)\right|\le \frac{C\epsilon h_n}{4D}\,|T^{q_n-j}(x)-x_n|\le \frac{C\epsilon h_n}{4CD} \frac{2D}{q_n}\le\epsilon/2
\]
and thus $|S_{q_n} f(x)-y|\le \epsilon$. It~remains to notice that
\begin{align*}
\liminf_{n\to\infty}\Leb(\Xi_n^y)&\ge \liminf_{n\to\infty}\frac{\epsilon h_n}{4CD}\frac{1}{\max\{C,D,E\}}\,|I_n|\\
&\ge \liminf_{n\to\infty} \frac{1}{4(\max\{C,D,E\})^2}\Leb(\Xi_n)>0.
\end{align*}
Thus we have proved that $y\in \Ess(T_f)$. Note that for every $n\in\N$ {we have}
\[
x\in \Bigl(x_n,x_n+\frac{1}{2\max\{C,D,E\} q_n}\Bigr).
\]
Hence,
\[
S_{q_n}f(x)-S_{q_n}f(x_n)\ge C^{-1}q_n |x-x_n|.
\]
Therefore, by~applying similar reasoning as to $y$, we~obtain that every
\[
z\in\Bigl[y,y+\frac{1}{2C^{-1}\max\{C,D,E\}}\Bigr]
\]
is an essential value of $T_f$. This finishes the proof of the proposition.
\end{proof}

\begin{thebibliography}{99}
\bibitem{berk_ergodicity_2024} P. Berk \& F. Trujillo - “On the ergodicity of infinite antisymmetric extensions of symmetric IETs”, J. Eur. Math. Soc. (JEMS) (2025), published online first
\bibitem{berk_ergodicity_2023} P. Berk, F. Trujillo \& C. Ulcigrai - “Ergodicity of explicit logarithmic cocycles over IETs”, 2023
\bibitem{boissy_dynamics_2009} C. Boissy \& E. Lanneau - “Dynamics and geometry of the Rauzy–Veech induction for quadratic differentials”, Ergodic Theory Dynam. Systems 29 (2009) no. 3, p. 767-816
\bibitem{conze_cocycles_2011} J.-P. Conze \& K. Frączek - “Cocycles over interval exchange transformations and multivalued Hamiltonian flows”, Adv. Math. 226 (2011) no. 5, p. 4373-4428
\bibitem{katok_interval_1980} A. Katok - “Interval exchange transformations and some special flows are not mixing”, Israel J. Math. 35 (1980) no. 4, p. 301-310
\bibitem{veech_interval_1978} W. A. Veech - “Interval exchange transformations”, J. Analyse Math. 33 (1978) no. 1, p. 222-272
\end{thebibliography}
\end{document}
